\section{Savremeni izazovi, dizajnerska retrospektiva i budući rad}
\label{sec:izazovi}
% TODO: Razraditi poglavlje tek nakon provere izvora; ne unositi neproverene činjenice.
Diskusija će razmotriti moguće drugačije projektne izbore bez pretpostavke
da bi oni nužno dali bolji standard. Predlozi budućeg rada ostaju ograničeni
na pitanja koja neposredno proizlaze iz pregleda i planiranih eksperimenata.

\subsection{Šta bismo danas možda uradili drugačije?}
% TODO: Istražiti svrhu AES-192, veći blok, key schedule i kompromis fleksibilnosti i jednostavnosti. Izneti različita mišljenja kao mišljenja potkrepljena argumentima; blogovi i Stack Exchange služe pronalaženju pitanja.

\subsection{Implementacije i AEAD kao polazište}
% TODO: Razmotriti side-channel-friendly dizajn, constant-time softver, akceleraciju, nonce disciplinu i autentifikovanu enkripciju kao primarni interfejs.

\subsection{Embedded, IoT i lightweight kontekst}
% TODO: Istražiti odnos AES-a i Ascon-a prema stvarnim ograničenjima i relevantnim NIST izdanjima. Ne pretpostaviti da novi algoritam univerzalno zamenjuje AES.

\subsection{Budući eksperimenti}
% TODO: Proširiti sadašnji AES-NI on/off eksperiment na različite CPU arhitekture; ARM extensions, embedded/FPGA, energiju i odvojeno AES-GCM merenje ostaviti izvan sadašnjeg obima.

\subsection{Dugoročna istraživačka pitanja}
% TODO: Predvideti model kvantnog napada, testiranje bočnih kanala i poređenje sa modernim simetričnim rešenjima; definisati šta bi takvo poređenje zahtevalo od podataka i opreme.
